\subsection{FUNCTIONS OF TWO ARGUMENTS}

A function of two arguments is a function whose domain is a set of ordered pairs (or, equivalently, a subset of a product). Let \(f\) be such a function; if \((x, y)\) is an element of the domain of \(f\) , the value \(f((x, y))\) of \(f\) at the element \((x, y)\) is generally denoted by \(f(x, y)\) . Let \(f\) be a function of two arguments, \(D\) its domain, and \(C\) its target. For each \(y\) let \(A_y\) be the set of all \(x\) such that \((x, y) \in D\) (that is, the section of \(\overline{D}^1\) at \(y\) (no. 1)). The mapping

\[
x \rightarrow f (x, y) \quad (x \in \mathrm{A} _ {y}, f (x, y) \in \mathrm{C})
\]

is called the partial mapping defined by \(f\) , with respect to the value \(y\) of the second argument, and is denoted by \(f(\bullet, y)\) , or \(f(., y)\) (or sometimes \(f_y\) ); we have \(f(\bullet, y)(x) = f(x, y)\) for all \((x, y) \in D\) . Similarly, for each \(x\) let \(B_x\) be the set of all \(y\) such that \((x, y) \in D\) . The mapping

\[
y \rightarrow f (x, y) \quad (y \in \mathbf {B} _ {x}, f (x, y) \in \mathbf {C})
\]

is called the partial mapping defined by \(f\) , with respect to the value \(x\) of the first argument, and is denoted by \(f(x, \bullet)\) , or \(f(x, \cdot)\) (or sometimes \(f_x\) ); we have \(f(x, \bullet)(y) = f(x, y)\) for all \((x, y) \in \mathbf{D}\) . If for every \(y\) (resp. \(x\) ) the partial mapping \(f(\bullet, y)\) (resp. \(f(x, \bullet)\) ) is a constant mapping, we say that \(f\) does not depend on its first (resp. second) argument; this means therefore that \(f(x, y) = f(x', y)\) whenever \((x, y)\) and \((x', y)\) belong to \(\mathbf{D}\) (resp. \(f(x, y) = f(x, y')\) whenever \((x, y)\) and \((x, y')\) belong to \(\mathbf{D}\) ). For each \(y\) belonging to the second projection of \(\mathbf{D}\) let \(g(y)\) denote the common value of the \(f(x, y)\) for \(x \in A_y\) ; the mapping \(y \to g(y)\) is then a mapping of \(\mathrm{pr}_2\mathbf{D}\) into \(\mathbf{C}\) such that

\[
g (y) = f (x, y) \quad \text {   for   all   } (x, y) \in D.
\]

Conversely, let \(g\) be a mapping of a set \(B\) into a set \(C\) , and let \(A\) be any set. Then the mapping \((x, y) \to g(y)\) of \(A \times B\) into \(C\) does not depend on its first argument.

Let \(u\) be a mapping of A into C and \(v\) a mapping of B into D. The mapping \(z \to (u(\mathrm{pr}_1z), v(\mathrm{pr}_2z))\) of \(A \times B\) into \(C \times D\) is called the (canonical) extension of \(u\) and \(v\) to the product sets, or simply the product of \(u\) and \(v\) (if there is no risk of confusion); its range is \(u\langle A\rangle \times v\langle B\rangle\) . It is denoted by \(u \times v\) . If \(u\) and \(v\) are injective (resp. surjective) then \(u \times v\) is injective (resp. surjective). If \(u\) and \(v\) are bijective, then \(u \times v\) is bijective, and its inverse mapping is \(\overline{u}^1 \times \overline{v}^1\) . If \(u'\) is a mapping of C into a set E and if \(v'\) is a mapping of D into a set F, we have

\[
\left(u ^ {\prime} \times v ^ {\prime}\right) \circ (u \times v) = \left(u ^ {\prime} \circ u\right) \times \left(v ^ {\prime} \circ v\right).
\]

If U, V are the graphs of u, v respectively, the graph W of \(u \times v\) is the set of ordered pairs \(((x, y), (z, t))\) of \((\mathbf{A} \times \mathbf{B}) \times (\mathbf{C} \times \mathbf{D})\) such that \((x, z) \in \mathbf{U}\) and \((y, t) \in \mathbf{V}\) ; the mapping \(((x, y), (z, t)) \to ((x, z), (y, t))\) puts W in one-to-one correspondence with the product \(U \times V\) (a subset of \((\mathbf{A} \times \mathbf{C}) \times (\mathbf{B} \times \mathbf{D})\) ) (cf.~§5, no. 5).

